% applications/railway-DE.tex

\chapter{Anwendungen für Eisenbahn-Alignments}
\label{chap:applications-railway}

Das Eisenbahnwesen ist die primäre Anwendungsdomäne, anhand derer diese
Arbeit die Fähigkeiten von AIM zeigt. Das Kapitel wendet die etablierte
Alignment-Semantik sowie die Zustands- und Laufzeitbegriffe auf praktische
Fragen des Eisenbahnwesens an.

\section{Anwendungskontext}

Die Verwendung von Eisenbahn-Alignments umfasst gekoppelte
Horizontalgeometrie, Höhenprofil, Überhöhung, Geschwindigkeit,
Zulässigkeitsbedingungen und Instandhaltungskontext. In AIM werden diese
als koordinierte Ebenen behandelt, die aus den vorangehenden Teilen
übernommen werden.

\section{Verwendung kanonischer Zustands- und Laufzeitbegriffe}

Eisenbahnanwendungen verwenden kanonische pose2/pose3- und Laufzeitdienste,
anstatt sie neu zu definieren. Betriebliche Größen werden mit den
etablierten Operatoren als abgeleitete Laufzeitgrößen ausgewertet.

Dies unterstützt Ingenieuraufgaben wie den Vergleich von Alignments, die
Bewertung der Übergangsqualität und eine realisierungsbewusste Auswertung
in bestehender Infrastruktur.

\section{Übergänge und gekoppeltes Verhalten in der Praxis}

Übergangsverhalten wird aus der Krümmungsentwicklung zusammen mit dem
Kontext von Überhöhung und Betriebsgeschwindigkeit interpretiert. Auf
Anwendungsebene gilt das Interesse nicht nur der geometrischen
Randwertübereinstimmung, sondern den Konsequenzen für Komfort,
Zulässigkeit und Lebenszyklusverhalten.

\section{Betriebliche und instandhalterische Konsequenzen}

In Anwendungsabläufen wirken Alignment-Entscheidungen mit Bedingungen an
Überhöhung, Übergangsqualität und Instandhaltbarkeit zusammen. AIM
unterstützt dies, indem es persistente Semantik von abgeleiteten
betrieblichen Größen trennt und eine laufzeit- und realisierungsbewusste
Auswertung ermöglicht.

\section{Interoperabilitätskontext}

Eisenbahnanwendungen erfordern häufig den Austausch mit BIM/GIS- und
Formatabläufen. In diesem Teil werden solche Standards als
Interoperabilitätsmechanismen behandelt, die der kanonischen
Ingenieursemantik nachgelagert sind.

\section{Verbindung zu AIM}

\begin{summary}
Eisenbahnanwendungen in AIM zeigen, wie Alignment-Semantik, Zustand,
Laufzeitauswertung und realisierungsbewusste Interpretation gemeinsam
praktische Ingenieurentscheidungen unterstützen.

Das Kapitel bleibt anwendungsbezogen und definiert keine kanonischen
Begriffe neu, deren Zuständigkeit bei vorgelagerten Teilen liegt.
\end{summary}
